% Umschlag für „Die polyzentrische Ordnung“, beide Bände.
%
% Diese Datei ist keine Pandoc-Vorlage. Das Satzskript schreibt ihr einen
% Kopf aus \def-Zeilen voran – Maße, Farbmodell, Texte – und setzt das Ganze
% mit LuaLaTeX in einem Lauf. Gezeichnet wird in Zoll, weil KDP in Zoll
% vermisst; der Ursprung liegt unten links auf dem Bogen, samt Beschnitt.
%
% Zwei Arten entstehen aus derselben Zeichnung:
%
% *Druck* (\UDruck = 1): der Umschlag des Taschenbuchs für KDP, ein einziger
% Bogen aus Rückseite, Rücken und Vorderseite, rundum mit 0,125 Zoll
% Beschnitt. Die Rückenbreite folgt aus der Seitenzahl des Innenteils und
% wird bei jedem Satz neu gerechnet; die Vorlage kennt nur ihr Ergebnis,
% \URuecken. Farben im CMYK-Modell, weil gedruckt wird.
%
% *E-Book* (\UDruck = 0): das Titelbild des E-Books, nur die Vorderseite,
% ohne Beschnitt, im Seitenverhältnis 1 : 1,6, das KDP für E-Book-Umschläge
% empfiehlt. Farben im RGB-Modell, weil Kindle kein CMYK kennt. Die
% Vorderseite ist deshalb in Anteilen der Höhe gebaut: Dieselbe Zeichnung
% trägt 8,5 und 8,8 Zoll.
%
% Was KDP vorschreibt und hier eingehalten wird: Text mindestens 0,125 Zoll
% innerhalb der Schnittkante (hier überall mehr), Rückentext nur bei mehr
% als 79 Seiten und mit mindestens 0,0625 Zoll Luft zu beiden Falzen – das
% Satzskript rechnet die Schriftgröße danach aus und lässt den Rückentext
% weg, wo er nicht passt –, ein weißes Feld von 2 × 1,2 Zoll für den
% Barcode unten an der Rückseite, 0,25 Zoll vom Rücken und von der unteren
% Schnittkante. In dieses Feld darf nichts gedruckt werden; KDP setzt den
% Barcode selbst hinein. Keine Transparenz – die hellen Kreise sind
% ausgemischte Volltöne, keine Deckkraftstufen –, weil KDP flache Dateien
% verlangt, und keine Linie unter 0,75 pt.

% fix-cm macht die Computer-Modern-Schriften stufenlos skalierbar. Sie
% kommen auf dem Umschlag gar nicht vor, aber LaTeX richtet für jede
% Schriftgröße auch die Mathematikschriften ein und meldete beim 35-Punkt-
% Titel sonst drei Ersetzungen, die niemand zu suchen braucht.
\RequirePackage{fix-cm}
\documentclass{article}

% Das Farbmodell kommt aus dem Kopf. \edef, weil LaTeX Paketoptionen nicht
% selbst expandiert.
\edef\UPaketoptionen{\noexpand\usepackage[\UFarbmodell]{xcolor}}
\UPaketoptionen

\usepackage{fontspec}
\usepackage[ngerman]{babel}
\usepackage{tikz}
\usepackage[
  paperwidth=\UBogenbreite in,
  paperheight=\UBogenhoehe in,
  margin=0pt,
  nohead, nofoot,
]{geometry}

% Ohne Objektströme steht die MediaBox im Klartext in der Datei. Das
% Satzskript liest sie dort nach und prüft, ob der Bogen das Maß hat, das
% KDP aus der Seitenzahl erwartet.
\pdfvariable objcompresslevel=0

\setmainfont{Alegreya}[
  Path = \USchriften,
  Extension = .otf,
  UprightFont = *-Regular,
  ItalicFont = *-Italic,
  BoldFont = *-Bold,
  BoldItalicFont = *-BoldItalic,
  SmallCapsFont = AlegreyaSC-Regular,
  Ligatures = TeX,
]

% --------------------------------------------------------------------------
% Farben
% --------------------------------------------------------------------------
% Jeder Band hat seinen Grund: der erste ein gedecktes Schieferblau, das die
% Akzentfarbe des Innenteils aufnimmt, der zweite ein dunkles Ziegelrot. Die
% Farbmenge des Grundes bleibt unter 240 Prozent, damit der Druck nicht
% absäuft.
%
% Jede Farbe steht zweimal, als CMYK für den Druck und als RGB für das
% Titelbild; xcolor nimmt die Angabe, die zum Farbmodell des Laufs passt.
% Umrechnen ließe xcolor sie auch, aber nach einer Formel, die dunkle Töne
% ins Schwarzblaue verschiebt. Die RGB-Werte sind stattdessen die, als die
% ein gewöhnliches Farbprofil die CMYK-Werte wiedergibt: So sieht das
% Titelbild aus wie das gedruckte Buch. Wer eine Farbe ändert, ändert beide.

\ifcase\UBand\relax
  \definecolor{grund}{cmyk/rgb}{0.60,0.40,0.40,0.50/0.259,0.322,0.333}% kein Band 0
\or
  \definecolor{grund}{cmyk/rgb}{0.75,0.50,0.35,0.45/0.188,0.290,0.357}
\or
  \definecolor{grund}{cmyk/rgb}{0.30,0.75,0.80,0.45/0.451,0.224,0.153}
\else
  \definecolor{grund}{cmyk/rgb}{0.60,0.40,0.40,0.50/0.259,0.322,0.333}
\fi
\definecolor{schrift}{cmyk/rgb}{0,0.02,0.06,0.02/0.976,0.957,0.914}
\definecolor{gold}{cmyk/rgb}{0.12,0.24,0.55,0.08/0.812,0.690,0.478}
\definecolor{papierweiss}{cmyk/rgb}{0,0,0,0/1,1,1}
\colorlet{kreislinie}{gold!58!grund}
\colorlet{kreisflaeche}{schrift!9!grund}
\colorlet{leisegold}{gold!80!grund}

\pagestyle{empty}
\setlength{\parindent}{0pt}
\setlength{\topskip}{0pt}
\frenchspacing

% --------------------------------------------------------------------------
% Maße, die TeX selbst ausrechnet
% --------------------------------------------------------------------------
% Der Satzspiegel der Vorder- und Rückseite: 0,58 Zoll Rand zu jeder
% Schnittkante. TikZ rechnet mit Zahlen in Zoll (\rand in der Zeichnung),
% die Kästen darunter mit TeX-Maßen; beide meinen dasselbe.
\newdimen\URand
\URand=0.58in
\newdimen\UTextbreite
\UTextbreite=\dimexpr\UBreite in - 2\URand\relax

% Ein Umbruch mit Flatterrand, der trotzdem trennt: \raggedright erlaubt dem
% Zeilenende unendlich viel Luft und trennt deshalb nie, und auf schmaler
% Spalte flattert das deutsche Kompositum dann ausladend.
\newcommand{\Flatter}{%
  \rightskip=0pt plus 2.5em
  \parfillskip=0pt plus 1fil
  \spaceskip=\fontdimen2\font minus \fontdimen4\font
}

\newsavebox\UTitelkasten
\newsavebox\UKlappe
\newsavebox\UWerk
\newdimen\UTitelgroesse
\newdimen\URest

\makeatletter
\begin{document}

% Der Titel steht in zwei Zeilen, umbrochen vor dem letzten Wort. Er bekommt
% 35 Punkt, es sei denn, seine längere Zeile wäre dann breiter als der
% Satzspiegel – dann so viel weniger, dass sie genau hineinpasst.
\sbox\UTitelkasten{\fontsize{35pt}{37pt}\selectfont\bfseries
  \begin{tabular}{@{}l@{}}\UTitelzeilen\end{tabular}}
\ifdim\wd\UTitelkasten>\UTextbreite
  \UTitelgroesse=\dimexpr 35pt*\UTextbreite/\wd\UTitelkasten\relax
\else
  \UTitelgroesse=35pt
\fi
\edef\UTitelschrift{\noexpand\fontsize{\the\UTitelgroesse}%
  {\the\dimexpr\UTitelgroesse*106/100\relax}\noexpand\selectfont\noexpand\bfseries}

\ifnum\UDruck=1
  % Der Umschlagtext und die Werkangabe darunter werden vorab gesetzt und
  % vermessen. Reicht der Platz zwischen dem oberen Rand und der Werkangabe
  % nicht, schreibt TeX eine Zeile ins Protokoll, die das Satzskript als
  % Warnung weitergibt – der Text liefe sonst stumm in die Werkangabe.
  % Die Farbe gehört in den Kasten selbst: Was vorab gesetzt ist, färbt der
  % Knoten, der es später aufnimmt, nicht mehr um.
  \sbox\UKlappe{\parbox[t]{\UTextbreite}{%
    \color{schrift}\fontsize{10.5pt}{14.6pt}\selectfont\Flatter
    \setlength{\parskip}{6pt}%
    \UKlappentext\par}}
  \sbox\UWerk{\parbox[b]{\UTextbreite}{%
    \color{gold}\rule{0.9in}{1pt}\par\vspace{7pt}%
    \fontsize{10pt}{13.5pt}\selectfont\UReihe\par}}
  \URest=\dimexpr \UHoehe in - 0.78in - \ht\UKlappe - \dp\UKlappe
    - 2.05in - \ht\UWerk - \dp\UWerk - 0.25in\relax
  \ifdim\URest<0pt
    \typeout{UMSCHLAG-UEBERLAUF: \strip@pt\dimexpr-\URest\relax}
  \fi
\fi
\makeatother

\noindent
\begin{tikzpicture}[x=1in, y=1in, line cap=round, line join=round]
  \useasboundingbox (0,0) rectangle (\UBogenbreite,\UBogenhoehe);

  % Der Grund läuft über den ganzen Bogen, Beschnitt eingeschlossen. Über
  % den Falz hinweg bleibt er dieselbe Fläche: KDP lässt dem Falz 0,0625
  % Zoll Spiel, und ein durchgehender Grund verzeiht es.
  \fill[grund] (0,0) rectangle (\UBogenbreite,\UBogenhoehe);

  % Die Vorderseite: Ursprung unten links an der Schnittkante.
  \pgfmathsetmacro{\fx}{\UBeschnitt + \UVorderseiteLinks}
  \pgfmathsetmacro{\fy}{\UBeschnitt}
  \pgfmathsetmacro{\w}{\UBreite}
  \pgfmathsetmacro{\h}{\UHoehe}
  \pgfmathsetmacro{\rand}{0.58}

  % ----------------------------------------------------------------------
  % Das Zeichen: Kreise mit je einem Mittelpunkt. Viele Zentren, keines in
  % der Mitte, keines um Längen das größte – das Motiv ist der Titel. Lage
  % in Anteilen der Seite, Halbmesser in Zoll; die Anordnung ist gesetzt und
  % nicht gewürfelt, damit jeder Satz dasselbe Bild zeigt. Zwei Kreise laufen
  % über die Schnittkante in den Beschnitt, auf dem Druckbogen bis an den
  % Rücken, nie darüber.
  % ----------------------------------------------------------------------
  \begin{scope}
    \clip (\fx,\fy-\UBeschnitt)
      rectangle (\fx+\w+\UBeschnitt,\fy+\h+\UBeschnitt);
    \foreach \ax/\ay/\r/\voll in {
        0.19/0.505/0.50/1,
        0.43/0.575/0.30/0,
        0.63/0.470/0.74/0,
        0.905/0.585/0.38/1,
        0.33/0.395/0.24/0,
        0.83/0.365/0.46/0,
        0.085/0.340/0.16/0,
        0.50/0.325/0.12/1,
        1.03/0.450/0.30/0,
        0.255/0.640/0.13/0%
      }{
      \pgfmathsetmacro{\mx}{\fx + \ax*\w}
      \pgfmathsetmacro{\my}{\fy + \ay*\h}
      \ifnum\voll=1
        \fill[kreisflaeche] (\mx,\my) circle[radius=\r];
      \fi
      \draw[kreislinie, line width=1.1pt] (\mx,\my) circle[radius=\r];
      \fill[gold] (\mx,\my) circle[radius=0.034];
    }
  \end{scope}

  % ----------------------------------------------------------------------
  % Vorderseite: Titel, Untertitel, Band, Verfasser
  % ----------------------------------------------------------------------
  \node[anchor=north west, inner sep=0pt, align=left, text=schrift,
        font=\UTitelschrift]
    at (\fx+\rand, \fy+0.905*\h) {\UTitelzeilen};

  % flush left ist der echte Flatterrand ohne Trennung; TikZs left trennt
  % und dehnt die Wortabstände, und ein getrennter Untertitel sieht nach
  % Versehen aus.
  \node[anchor=north west, inner sep=0pt, text width=\UTextbreite,
        align=flush left, text=gold]
    at (\fx+\rand, \fy+0.735*\h) {%
      \fontsize{14.5pt}{18.5pt}\selectfont\itshape\UUntertitel\par};

  \node[anchor=south west, inner sep=0pt, text width=\UTextbreite,
        align=flush left]
    at (\fx+\rand, \fy+0.155*\h) {%
      {\color{gold}\fontsize{12.5pt}{14pt}\selectfont\scshape
        \addfontfeatures{LetterSpace=6}\UBandzahl\par}%
      \vspace{3pt}%
      {\color{schrift}\fontsize{21pt}{24pt}\selectfont\UBandtitel\par}};

  \node[anchor=south west, inner sep=0pt, text width=\UTextbreite,
        align=flush left, text=schrift]
    at (\fx+\rand, \fy+0.062*\h) {%
      \fontsize{14pt}{16pt}\selectfont\scshape
      \addfontfeatures{LetterSpace=4}\UAutor\par};

  % ----------------------------------------------------------------------
  % Rücken und Rückseite gibt es nur auf dem Druckbogen.
  % ----------------------------------------------------------------------
  \ifnum\UDruck=1

    % Rücken: Verfasser unten, Titel in der Mitte, Band oben, von unten
    % nach oben zu lesen, wie es in deutschen Büchern üblich ist.
    \ifnum\URueckentext=1
      \pgfmathsetmacro{\rx}{\UBeschnitt + \UBreite + \URuecken/2}
      \node[rotate=90, anchor=mid west, inner sep=0pt, text=schrift]
        at (\rx, \UBeschnitt+0.55) {%
          \fontsize{\URueckengroesse pt}{\URueckengroesse pt}\selectfont\scshape
          \addfontfeatures{LetterSpace=4}\UAutor};
      \node[rotate=90, anchor=mid, inner sep=0pt, text=schrift]
        at (\rx, \UBeschnitt+\UHoehe/2) {%
          \fontsize{\URueckengroesse pt}{\URueckengroesse pt}\selectfont\bfseries
          \UTitel};
      \node[rotate=90, anchor=mid east, inner sep=0pt, text=gold]
        at (\rx, \UBeschnitt+\UHoehe-0.55) {%
          \fontsize{\URueckengroesse pt}{\URueckengroesse pt}\selectfont\scshape
          \addfontfeatures{LetterSpace=4}\URueckenband};
    \fi

    % Rückseite: Ursprung unten links an der Schnittkante.
    \pgfmathsetmacro{\bx}{\UBeschnitt}
    \pgfmathsetmacro{\by}{\UBeschnitt}

    \node[anchor=north west, inner sep=0pt, text=schrift]
      at (\bx+\rand, \by+\h-0.78) {\usebox\UKlappe};

    % Die Werkangabe: in welchem Band man liest und welcher dazugehört.
    \node[anchor=south west, inner sep=0pt, text=gold]
      at (\bx+\rand, \by+2.05) {\usebox\UWerk};

    % Nachweis links unten, neben dem Barcodefeld, in festen Zeilen: Die
    % Adresse des Verzeichnisses darf nicht am Bindestrich umbrechen.
    \node[anchor=south west, inner sep=0pt, align=flush left, text=leisegold,
          font=\fontsize{8.5pt}{11pt}\selectfont]
      at (\bx+\rand, \by+0.30) {\UNachweis};

    % Das Barcodefeld: 2 × 1,2 Zoll, 0,25 Zoll vom Rücken und von der
    % unteren Schnittkante, weiß und leer.
    \fill[papierweiss]
      (\bx+\w-0.25-2.0, \by+0.25) rectangle (\bx+\w-0.25, \by+0.25+1.2);

  \fi
\end{tikzpicture}
\end{document}
